% Deliberately neutral demo prose, not translated from a published paper.
\newcommand{\SampleOpening}{%
 \section*{引言：示例正文与引文}
 这是一段用于检查双栏阅读节奏的中性示例文字。段落中的术语、缩写和标点需要在实际项目中与源文保持一致。中文译稿应呈现原论文的完整论述，而不是把多段证据压缩成一段摘要。这里的引文只链接到本文件末尾明确标为占位的参考文献\SampleCiteOne。\par
 若原句仅表示相关、可能或提示，译文也应保留相同的证据强度。数字、单位、正负号、上下标和范围端点应逐项核对。原文存在内部矛盾时，应分别保留相应位置的陈述，再用简短译注说明；不能根据常识把不同的数字静默统一。\par
 \subsection*{内容层级与跨页衔接}
 此小标题演示二级层级。实际论文可能采用不同的小标题密度，模板不会替代对原文结构的识别。双栏文章需按左栏、右栏和跨页续句的顺序读取。若段落在页末中断，应在完成整句核对后接入下一页，避免把图中文字误当正文，或重复翻译同一段。\par
 图表具有独立的信息范围。正文引用到某幅图，并不意味着其图注可以省略。图注中的每个分面、比较对象、时间点、样本量及统计说明都应对应源文。英文图内标签可以保留，并在图旁或文末提供中文释义；实际采用哪种安排，需要依据图形尺寸和译文长度判断\SampleCiteTwo。\par
 \subsection*{来源与可维护数据}
 本段继续演示正文密度。可编辑文本和图注数据应与最终版面分开维护，以便修订时准确定位变化。若只调整版式，应能证明正文、图注和书目内容没有被意外改写。更换字体或纸张后，须重新检查断行、图注与图片的相对位置以及页脚可读性。\par
 文内引用与书目之间应有明确对应关系。这里的两条演示书目没有真实DOI，也没有虚构的实验信息。使用模板时可增删条目，但必须同时更新编号、链接目标和引用位置。页码只表示中文译稿自己的分页，不用于冒充原出版页码。\par
 \subsection*{结果叙述的演示占位}
 本节的文字仅用于呈现长段落在双栏页面中的流动，不报告真实结果。实际翻译应保留观察对象、实验条件和结局指标之间的关系。例如，原文分别描述体外实验与动物实验时，不能把两类结果合并为同一层面的结论；研究使用的对照、时间顺序和分析范围也必须保留。图表中的比较范围不应扩大为对整个人群或全部疾病阶段的判断。\par
 论述中的转折与限制同样属于正文。若研究先提出一种解释，再讨论另一种可能，译文应沿用这一推进顺序。作者明确写出的不确定性应能在中文中被识别；如果尚未确定某个现象的方向或机制，就不能改成已经建立的因果关系。对于术语近似但研究层级不同的概念，可在首次出现时保留缩写，后文保持统一。\par
 \subsection*{方法信息与可复核性}
 方法部分往往包含密集数字和专有名词。翻译时须区分浓度、总量、频率和处理时长，不能因为中文句式调整而改变参数的归属。试剂名称、软件版本和数据访问标识通常保留原样；上下限、排除规则和显著性标准则需要逐字符核对。这里没有提供任何真实实验参数，避免把模板占位误认为可以直接执行的方法。\par
 表格与补充材料的引用也需要保留原层级。正文对补充图的指向并不意味着模板中已经附有该补充图；若源文件只包含主文，应如实标注所交付的范围。图表编号、附录顺序和参考文献数量都应由目标文章决定，本示例中的结构不构成固定配额，也不要求所有论文采用同样的章节。\par
}
\newcommand{\SampleClosing}{%
 \section*{讨论：版式适配与范围}
 上方矢量图是模板原创的排版占位图，不来自任何论文。图形与完整演示图注被放在同一页，便于读者直接核对分面说明。真实项目应使用经过授权或用户提供的原图素材，并保留图注署名与许可信息。尺寸不足时，可适度缩放或调整栏宽，不以删除图注解决空间问题。\par
 模板仅记录一次已核验文章的可观察设计特征。即使文章出自同一期刊，不同年代、研究论文与综述、校正清样与正式版本也可能采用不同的页眉和引文形式。应以用户提供的实际PDF为准，不能仅凭期刊名称推断配色、字号或引用方式\SampleCiteOne。\par
 \subsection*{演示检查项}
 本样板用于展示可编译性、中文字形、标题层级、跨栏图注、页眉页脚和内部链接。图中没有实际数据，因此不应被当作科学插图引用。替换内容后，还需检查所有页的渲染结果，确认无截字、碰撞、异常空白和遗漏段落。\par
 \SampleBibliography
 \SampleSourceInfo
}
